\chapter{ТЕОРЕТИЧЕСКИЕ ОСНОВЫ И МЕТОДЫ ОПТИМИЗАЦИИ МОДЕЛЕЙ ДЛЯ КЛИЕНТСКОЙ СРЕДЫ}
\label{ch:теоретические_основы_и_методы_оптимизации_моделей_для_клиентской_среды}

\section{Эволюция методов обработки естественного языка и архитектура трансформеров}
\label{sec:эволюция_методов_обработки_естественного_языка_и_архитектура_трансформеров}

Обработка естественного языка (Natural Language Processing, NLP) прошла путь от эвристических правил и статистических моделей до глубоких нейронных архитектур, формирующих богатые контекстуальные представления текста. Долгое время стандартом де-факто служили подходы, требующие явного конструирования признаков вручную: векторное представление мешка слов (Bag-of-Words), $n$-граммные модели и классификаторы на опорных векторах (SVM). Принципиальный изъян такого семейства методов – неспособность учитывать зависимости между словами, разнесёнными в тексте на произвольное расстояние. Для морфологически богатых языков со свободным порядком слов этот недостаток оборачивается существенным падением точности при решении задач анализа тональности.

Рекуррентные нейронные сети (RNN) вместе со своими модификациями – LSTM (долгосрочная краткосрочная память) и GRU (управляемые рекуррентные блоки) – ознаменовали переход к векторным представлениям, извлекаемым из последовательностей в ходе обучения. В основе всех трёх архитектур лежит пошаговое обновление скрытого состояния: $h_t = f(W_h h_{t-1} + W_x x_t + b)$, что теоретически позволяет сохранять контекст на протяжении сколь угодно длинной цепочки токенов. Оборотная сторона такой пошаговости – принципиальная невозможность параллельного исполнения на GPU и склонность к экспоненциальному затуханию градиентов при обратном распространении ошибки сквозь время (BPTT). В контексте браузерного развёртывания этот изъян проявляется особенно остро: синхронный характер инференса блокирует поток исполнения и делает LSTM-архитектуры малопригодными для интерактивных веб-приложений, где требуется отзывчивый пользовательский интерфейс.

В работе~\cite{vaswani2017attention}, опубликованной в 2017 году, от рекуррентности отказались полностью: её место занял механизм самовнимания (Self-Attention), вычисляющий попарные зависимости между всеми позициями входной последовательности за один шаг без какой-либо последовательной обработки. Операция масштабированного скалярного произведения (Scaled Dot-Product Attention) записывается как~\eqref{eq:scaled_dot_product_attention}
\begin{equation}
    \mathrm{Attention}(Q, K, V) = \mathrm{softmax}\!\left(\frac{QK^\top}{\sqrt{d_k}}\right)V,
    \label{eq:scaled_dot_product_attention}
\end{equation}
где $\mathrm{Attention}(Q, K, V)$ – выходная матрица контекстно взвешенных представлений; $\mathrm{softmax}(\cdot)$ – преобразование scores в вероятностное распределение; $Q$, $K$, $V$ – матрицы запросов, ключей и значений соответственно; $d_k$ – размерность ключевого пространства; деление на $\sqrt{d_k}$ удерживает аргументы $\mathrm{softmax}$ в ненасыщенной области. Механизм многоголового внимания (Multi-Head Attention) параллельно задействует несколько независимых подпространств признаков: $\mathrm{MultiHead}(Q,K,V)=\mathrm{Concat}(\mathrm{head}_1,\ldots,\mathrm{head}_h)W^O$. Принципиальное преимущество архитектуры состоит в том, что путь между любой парой позиций имеет длину $O(1)$, тогда как в рекуррентных сетях та же величина растёт как $O(N)$ – именно это делает трансформер надёжным инструментом для улавливания отрицаний и дальних грамматических зависимостей в задачах анализа тональности.

Архитектура трансформера послужила основой для создания модели BERT (Bidirectional Encoder Representations from Transformers)~\cite{devlin2019bert}, в которой задействован двунаправленный стек энкодеров с функцией активации GELU, остаточными соединениями и нормализацией по слоям: $\mathrm{Output} = \mathrm{LayerNorm}(x + \mathrm{Sublayer}(x))$. Двухэтапная схема предобучения~– на задаче восстановления замаскированных токенов (MLM) и на задаче определения смежности предложений (NSP)~– позволяет получить универсальные контекстуальные представления, которые затем адаптируются к классификации тональности посредством fine-tuning. Вместе с тем 110 миллионов параметров базовой конфигурации BERT-base обусловливают высокую стоимость загрузки весов и значительный пиковый расход памяти, что выводит модель за рамки допустимого для интерактивных клиентских приложений.

Данный барьер преодолевается применением архитектуры DistilBERT~\cite{sanh2019distilbert}, основанной на дистилляции знаний: сокращение числа слоёв энкодера с 12 до 6 при сохранении ширины скрытого состояния снижает число параметров на 40\% и ускоряет вывод модели на 60\% при сохранении 97\% точности. Компактные архитектуры семейства DistilBERT, конвертированные в формат ONNX и квантованные до INT8, составляют основу систем клиентского анализа тональности, реализуемых средствами Transformers.js и ONNX Runtime Web.

Параллельно с развитием моноязычных моделей сформировалось направление мультиязычных трансформеров, ориентированных на совместное представление текстов множества языков в едином векторном пространстве~\cite{conneau2020unsupervised}. Модель mBERT обучена на текстах 104 языков с общим словарём из 119~547 токенов: кросс-языковое обобщение достигается ценой значительного роста матрицы эмбеддингов~– порядка 91 миллиона параметров,~– не затрагиваемого процедурой дистилляции. Архитектура XLM-RoBERTa~\cite{conneau2020unsupervised}, предобученная на корпусе CC-100 объёмом 2,5 терабайта без задачи NSP, превосходит mBERT на большинстве кросс-языковых бенчмарков; дообученная версия cardiffnlp/twitter-xlm-roberta-base-sentiment ориентирована на трёхклассовую классификацию тональности. Вместе с тем исходный размер весов модели в формате FP32 превышает 1~ГБ, что обусловливает обязательность квантования перед развёртыванием через Transformers.js~\cite{xu2024distillation}.

Вычислительная сложность механизма самовнимания квадратично зависит от длины входной последовательности~– $O(N^2 \cdot d)$,~– что определяет нелинейный рост латентности инференса при увеличении числа токенов. Морфологически богатые языки, в частности русский и арабский, генерируют существенно более длинные последовательности при токенизации алгоритмом WordPiece, нежели аналитические языки типа английского~\cite{rust2021multilingual}: агглютинативная структура арабского письма приводит к фрагментации одного графического слова на 3–4 субтокена, резко увеличивая $N$ и нагрузку на вычислительный поток Web Worker. Разреженные и линейные схемы внимания~\cite{treviso2022efficient} предлагают решение данной проблемы, однако их поддержка в экосистеме Transformers.js и ONNX Runtime Web носит ограниченный характер; на практике применяется ограничение длины контекста в сочетании с квантованием и дистилляцией~\cite{gholami2021quantization}. Таким образом, выбор архитектуры для клиентского анализа тональности обусловлен совокупностью ограничений среды исполнения браузера, рассматриваемых в последующих разделах.

\section{Специфика мультиязычных моделей} \label{sec:специфика_мультиязычных_моделей}

Развитие мультиязычных трансформерных моделей обусловлено практической необходимостью обработки типологически разнородных языков в рамках единой архитектуры без поддержки отдельных моноязычных систем для каждого языка. Совместное предобучение на конкатенированных корпусах вынуждает модель формировать выровненные межъязыковые представления в общем векторном пространстве~\cite{conneau2020unsupervised}, открывая возможность дообучения на размеченных данных одного языка с последующим переносом на остальные~– свойство, принципиально важное при дефиците разметки для морфологически сложных языков, таких как русский и арабский.

Словарный состав мультиязычных моделей формируется алгоритмом WordPiece на основе объединённого многоязычного корпуса. Модель mBERT оперирует словарём из 119~547 токенов против 30~522 у BERT-base-uncased, а дистиллированный вариант multilingual DistilBERT сохраняет тот же словарь при сокращённом числе слоёв энкодера~\cite{sanh2019distilbert}. Матрица эмбеддингов размерности $V \times d$, где $d = 768$ для base-конфигурации, не подвергается сжатию в процессе дистилляции и содержит порядка 91 миллиона параметров: четырёхкратный рост словаря относительно моноязычного аналога влечёт линейное увеличение доли параметров, обязательно загружаемых в память при инициализации, но не участвующих в вычислениях слоёв внимания. В браузерной среде данное обстоятельство проявляется в двух взаимосвязанных эффектах: увеличении времени «холодного старта» вследствие роста размера бинарного файла ONNX и повышенном давлении на кучу WebAssembly, где векторы невостребованных языков занимают выделенный объём памяти~\cite{pfeiffer2022lifting}.

Скорость токенизации напрямую зависит от морфологического устройства языка, поскольку именно оно определяет длину входной последовательности $N$, а через квадратичную сложность $O(N^2 \cdot d)$ нелинейно влияет на задержку инференса~\cite{rust2021multilingual}. Аналитический строй английского языка задаёт нижнюю границу: небольшое число морфем на слово минимизирует отношение токенов к словам и удерживает длину последовательности на низком уровне. Русский язык с его развитой флективной системой вынуждает токенизатор WordPiece расщеплять словоформы по границе основы и окончания, из-за чего среднее количество субтокенов на предложение заметно возрастает и вычислительная нагрузка на поток Web Worker увеличивается пропорционально. Наиболее острую проблему для клиентского инференса создаёт арабская система письма: артикли, предлоги и местоимённые клитики присоединяются к опорному слову без пробела, что влечёт агрессивную сегментацию~– одно орфографическое слово нередко распадается на 3–4 субтокена, резко увеличивая $N$ и провоцируя квадратичный рост расхода памяти~\cite{farhan2024arabic}.

Обучение XLM-RoBERTa~\cite{conneau2020unsupervised} на корпусе CC-100 отличается от классической схемы BERT двумя ключевыми решениями: отказом от задачи NSP и применением динамического маскирования. Совокупность этих изменений даёт более робастные контекстуальные представления, чем у mBERT,~– преимущество прослеживается на большинстве кросс-языковых тестов и наиболее ощутимо на коротких жанрах~– отзывах и постах в социальных сетях,~– где тональность выражается через идиомы и разговорную лексику. Вместе с тем полноразмерный XLM-RoBERTa-base в формате FP32 занимает свыше 1~ГБ, и это делает квантование до INT8 обязательным условием развёртывания через ONNX Runtime Web~\cite{gholami2021quantization}. Сопоставление кандидатных архитектур для клиентского анализа тональности обнаруживает характерный компромисс между точностью и потреблением ресурсов~\cite{qin2024multilingual}: Xenova/bert-base-multilingual-uncased-sentiment держит умеренный расход памяти, однако на морфологически сложных языках уступает моделям семейства RoBERTa; Xenova/distilbert-base-multilingual-cased-sentiments-student показывает лучшее соотношение задержки и точности в браузерной среде; дообученный вариант ayayousef/distilbert-multilingual-sentiment-finetuned частично восполняет просадку точности, типичную для дистилляции в многоязычных сценариях~\cite{xu2024distillation}.

Самостоятельную сложность представляет дисбаланс языков в предобучающем корпусе. Эффект, получивший название «проклятия многоязычности» (curse of multilinguality)~\cite{pfeiffer2022lifting}, состоит в следующем: наращивание числа поддерживаемых языков при неизменной ёмкости модели неизбежно снижает качество представлений для каждого конкретного языка относительно моноязычной модели сопоставимого масштаба. Арабский и русский языки находятся на периферии частотного распределения обучающих данных, что влечёт два взаимосвязанных следствия: редкие словоформы непропорционально часто отображаются в токен [UNK], а семантически смежные единицы располагаются в пространстве эмбеддингов с пониженной точностью. Практическим итогом оказывается систематическое занижение метрики $F_1$ на этих языках по сравнению со специализированными моноязычными моделями~\cite{rust2021multilingual}~– фактор, который необходимо учитывать при анализе результатов экспериментальной валидации.

\section{Методы оптимизации для граничных вычислений посредством дистилляции знаний и квантования}
\label{sec:методы_оптимизации_для_граничных_вычислений_посредством_дистилляции_знаний_и_квантования}

Браузерная среда исполнения накладывает на развёртывание трансформерных моделей жёсткие ресурсные ограничения, принципиально отличающие её от серверной инфраструктуры: ограниченный объём кучи WebAssembly, однопоточная модель исполнения JavaScript и отсутствие прямого доступа к аппаратным ускорителям в совокупности исключают использование полноразмерных архитектур уровня BERT-base без предварительной оптимизации. Среди методов оптимизации для граничных вычислений (Edge Computing) наибольшее распространение получили дистилляция знаний (Knowledge Distillation) и квантование (Quantization)~\cite{menghani2023efficient}: оба метода преследуют единую цель снижения вычислительной ёмкости при минимальной деградации точности, реализуя её на разных уровнях представления модели.

Дистилляция знаний реализует парадигму «учитель~– ученик» (Teacher-Student): компактная модель-«ученик» обучается аппроксимировать не только целевые метки, но и полное вероятностное распределение выходного слоя крупной модели-«учителя»~\cite{xu2024distillation}. Ненулевые вероятности, присваиваемые «учителем» неверным классам, отражают структурные семантические зависимости~– «тёмные знания» (dark knowledge),~– усвоение которых позволяет «ученику» достигать обобщающей способности, несоразмерной его параметрической ёмкости. Для управления степенью сглаживания распределения вводится гиперпараметр температуры $T$ в функцию softmax, определяемый выражением~\eqref{eq:soft_targets}
\begin{equation}
    q_i = \frac{\exp(z_i / T)}{\sum_j \exp(z_j / T)},
    \label{eq:soft_targets}
\end{equation}
где $q_i$~– сглаженная вероятность $i$-го класса, $z_i$~– логиты модели-«учителя», $T$~– температура сглаживания, при $T > 1$ формируются мягкие метки (soft targets), акцентирующие межклассовые сходства. Суммарная функция потерь при обучении «ученика» объединяет три компонента, формализованных выражением~\eqref{eq:distillation_loss}
\begin{equation}
    \mathcal{L} = \alpha \mathcal{L}_{\mathrm{CE}} + \beta \mathcal{L}_{\mathrm{distill}} + \gamma \mathcal{L}_{\mathrm{cos}},
    \label{eq:distillation_loss}
\end{equation}
где где $\mathcal{L}$~– суммарная функция потерь, $\mathcal{L}_{\mathrm{CE}}$~– стандартная кросс-энтропия на размеченных данных, $\mathcal{L}_{\mathrm{distill}}$~– расхождение Кульбака~– Лейблера между распределениями «учителя» и «ученика» при температуре $T$, $\mathcal{L}_{\mathrm{cos}}$~– косинусное расстояние между скрытыми состояниями (hidden states) обеих моделей, обеспечивающее согласование внутренних представлений~\cite{sanh2019distilbert}; весовые коэффициенты $\alpha$, $\beta$, $\gamma$ подбираются эмпирически.

Применительно к архитектуре DistilBERT структурная компрессия достигается сокращением числа слоёв энкодера с 12 до 6 при сохранении $d = 768$ и $h = 12$; инициализация весов «ученика» производится выборкой каждого второго слоя «учителя», что обеспечивает «тёплый старт» и ускоряет сходимость. Удаление компонента NSP и соответствующих эмбеддингов типов токенов снижает параметрическую избыточность без ущерба для задач классификации на уровне предложения: совокупный эффект выражается в сокращении параметров на 40\% и ускорении вывода модели на 60\% при сохранении 97\% точности базовой модели~\cite{sanh2019distilbert}. В контексте браузерного инференса через Transformers.js и ONNX Runtime Web данная компрессия обеспечивает приемлемую латентность на устройствах среднего ценового сегмента~\cite{wang2024anatomizing}. В многоязычной среде передача кросс-языковых представлений от «учителя» к «ученику» осуществляется через минимизацию $\mathcal{L}_{\mathrm{cos}}$~\cite{xu2024distillation}: инвариантность относительно языка сохраняется, что позволяет модели~– в частности, Xenova/distilbert-base-multilingual-cased-sentiments-student~– удерживать устойчивые показатели точности на морфологически сложных языках при потреблении ресурсов, сопоставимом с моноязычными архитектурами. Матрица эмбеддингов $V \times d$ при $V = 119\,547$ в ходе дистилляции не изменяется и доминирует в профиле памяти~– этот дисбаланс устраняется на последующем этапе квантования.

Квантование сводится к отображению непрерывного диапазона вещественных значений на конечное множество дискретных уровней~\cite{gholami2021quantization}. В библиотеках ONNX Runtime и Transformers.js наибольшее распространение получила схема линейного квантования в формат INT8, задаваемая преобразованием~\eqref{eq:quantization}

\begin{equation}
    q = \mathrm{clamp}\!\left(\mathrm{round}\!\left(\frac{r}{S} + Z\right),\, q_{\min},\, q_{\max}\right),
    \label{eq:quantization}
\end{equation}
где $q$~– квантованное целочисленное значение, $r$~– исходное вещественное значение, $S$~– коэффициент масштабирования, $Z$~– целочисленное смещение, обеспечивающее точное представление нуля при работе с padding-заполнителями; $q_{\min}$ и $q_{\max}$~– границы допустимого диапазона, для INT8 стандартно равные $[-128,\, 127]$. Обратное преобразование $\hat{r} = S \cdot (q - Z)$ восстанавливает приближённое вещественное значение; возникающий шум квантования $\|\hat{r} - r\|$ подавляется внутренней избыточностью параметрического пространства языковых моделей при условии адекватной калибровки~\cite{dettmers2022llmint8}.

Методы квантования после обучения (Post-Training Quantization) принято делить на два класса~\cite{gholami2021quantization}. Динамическое квантование заблаговременно переводит веса слоёв в формат INT8, откладывая вычисление параметров масштабирования активаций до момента каждого прямого прохода: такой подход избавляет от необходимости в калибровочной выборке и сохраняет работоспособность при широкой вариативности входов~– свойство особенно ценное в мультиязычных сценариях, где статистика активаций существенно расходится от языка к языку. Статическое квантование идёт по иному пути: коэффициенты масштабирования $S$ и сдвиги $Z$ для активаций определяются заранее на представительной калибровочной выборке, что полностью снимает накладные расходы на сбор статистики во время инференса, однако оставляет риск деградации точности на данных, выходящих за границы калибровочного распределения.

Замена формата FP32 на INT8 вчетверо уменьшает объём памяти, занимаемой весами модели, и кардинально меняет картину исполнения в среде WebAssembly: компактные INT8-матрицы лучше вписываются в кэш-иерархию процессора, сокращая число промахов кэша (cache misses) и тем самым ускоряя операции матричного умножения (GEMM)~\cite{jia2024empowering}. Дополнительное ускорение обеспечивают инструкции WASM SIMD, позволяющие обрабатывать за один такт 16 восьмибитных значений вместо четырёх 32-битных~– теоретический потолок прироста производительности при переходе на INT8 составляет 400\%~\cite{jia2024empowering}. В реализации ONNX Runtime Web для предотвращения целочисленного переполнения применяются специализированные микроядра, которые накапливают частичные суммы в 32-битных аккумуляторах.

Совместное применение дистилляции знаний и динамического квантования INT8 формирует двухуровневую стратегию оптимизации, при которой дистилляция устраняет структурную избыточность на уровне архитектуры, а квантование снижает ресурсоёмкость на уровне представления параметров~\cite{tang2024transformer}. Конвертация ayayousef/distilbert-multilingual-sentiment-finetuned в формат ONNX с применением динамического квантования весов позволяет достичь латентности менее 200~мс на одно предложение даже на мобильных устройствах при сохранении способности корректно идентифицировать тональность на всех трёх целевых языках~\cite{wang2024anatomizing}; квантованный граф исполняется внутри Web Worker, изолируя тензорные операции от цикла событий браузера. Вместе с тем оба метода оптимизации сопряжены с компромиссами, требующими экспериментальной валидации: дистилляция влечёт деградацию точности на морфологически сложных языках вследствие эффекта «проклятия многоязычности», тогда как квантование вносит шум, степень влияния которого варьируется в зависимости от распределения активаций~\cite{xu2022survey}.

\section{Формат ONNX и графовые вычисления} \label{sec:формат_onnx_и_графовые_вычисления}

Браузерное окружение не располагает ни нативной Python-инфраструктурой, ни прямым доступом к CUDA-интерфейсам, что исключает применение экосистем PyTorch и TensorFlow в том виде, в каком они существуют на сервере. Формат ONNX (Open Neural Network Exchange) снимает это ограничение: он служит аппаратно-независимым промежуточным представлением вычислительных графов и открывает путь к исполнению моделей через ONNX Runtime Web~\cite{onnxruntime2024}. Именно исчерпывающее покрытие операторов, используемых в архитектурах BERT, DistilBERT и XLM-RoBERTa, превратило ONNX в де-факто стандарт для клиентского инференса в веб-приложениях~\cite{xu2022survey}.

В формате ONNX вычислительный граф модели фиксируется как детерминированный ориентированный ациклический граф (DAG) и сериализуется посредством бинарного протокола Protocol Buffers: компактное представление данных сокращает время загрузки и десериализации модели в браузере. В отличие от динамического режима обучения (eager execution), статически закреплённая топология графа допускает применение оптимизационных преобразований ещё до первого прямого прохода, что делает задержку инференса предсказуемой~\cite{wang2024anatomizing}. Механизм динамических осей (Dynamic Axes), реализованный через символьные переменные \texttt{batch\_size} и \texttt{sequence\_length}, устраняет жёсткую привязку к конкретной длине входной последовательности: граф остаётся пригодным для исполнения без перекомпиляции при любом числе токенов~--- что особенно важно для NLP-задач, где длина токенизированного текста определяется морфологией обрабатываемого языка.

До передачи графа в среду исполнения над ним последовательно выполняются преобразования нескольких уровней~\cite{onnxruntime2024}. Первый уровень нацелен на устранение семантической избыточности, привнесённой при экспорте: свёртка констант (Constant Folding) заменяет статические подграфы их заранее вычисленными значениями, а удаление мёртвого кода (Dead Code Elimination) исключает операторы, результаты которых нигде не используются. Для трансформерных архитектур ключевую роль играет слияние операторов (Node Fusion): отдельные операции Multi-Head Attention, LayerNorm и GELU сворачиваются в единые составные ядра (fused kernels), что удерживает промежуточные результаты в кэше процессора и существенно разгружает шину памяти~\cite{jia2024empowering}. На третьем уровне применяется оптимизация размещения тензоров в памяти (Layout Optimization), обеспечивающая последовательный доступ к данным и тем самым повышающая эффективность векторизации операций матричного умножения.

Конвертация моделей~– Xenova/bert-base-multilingual-uncased-sentiment, Xenova/distilbert-base-multilingual-cased-sentiments-student, cardiffnlp/twitter-xlm-roberta-base-sentiment~– в формат ONNX осуществляется инструментарием optimum библиотеки Hugging Face с сохранением динамических осей~\cite{tang2024transformer}. Применение динамического квантования весов до INT8 непосредственно при экспорте сокращает размер бинарного файла в четыре раза относительно версии FP32, что линейно снижает время «холодного старта» (cold start)~– критическую метрику для клиентских веб-приложений, где задержка инициализации непосредственно влияет на воспринимаемое качество пользовательского опыта~\cite{gholami2021quantization}.

Исполнение оптимизированного графа обеспечивается ONNX Runtime Web~– кроссплатформенным движком с модульной системой провайдеров исполнения (Execution Providers)~\cite{onnxruntime2024}. В рамках разрабатываемой системы в качестве основного бэкенда задействуется провайдер WASM с поддержкой инструкций SIMD: данный выбор обусловлен отсутствием накладных расходов на копирование данных между оперативной и видеопамятью, широкой поддержкой WASM SIMD в браузерах на базе Chromium и Firefox, а также эффективностью целочисленных вычислений INT8 на CPU при малом размере батча, характерном для интерактивного анализа тональности~\cite{wang2024anatomizing}.

Для предотвращения блокировки основного потока интерфейса (UI blocking) сессия ONNX Runtime Web инициализируется и исполняется внутри Web Worker, изолируя тензорные вычисления от цикла событий браузера~\cite{jia2024empowering}; входные тензоры передаются через механизм Transferable Objects без копирования буферов, минимизируя накладные расходы на межпоточную коммуникацию. Высокоуровневая обёртка Transformers.js автоматизирует токенизацию, формирование тензоров input\_ids, attention\_mask и token\_type\_ids, запуск сессии ONNX Runtime и постобработку логитов, предоставляя унифицированный API для инференса моделей различных архитектур~\cite{ruan2024webllm}; непосредственная вычислительная нагрузка при этом целиком ложится на механизмы ONNX Runtime, тогда как Transformers.js выполняет роль оркестратора, абстрагирующего от аппаратных особенностей устройства пользователя. Совокупность описанных механизмов~– статический граф с динамическими осями, многоуровневая оптимизация операторов, квантование весов до INT8 и исполнение через провайдер WASM SIMD в изолированном потоке Web Worker~– формирует целостный технологический стек, обеспечивающий практическую реализуемость клиентского вывода мультиязычных трансформеров в браузерной среде~\cite{xu2022survey}.